plex@@PAGE_MARK@@ analytic variant of the Riemann--Roch theorem presented here, or suitable variants of the Atiyah--Singer index theorem. Some indications in this direction will be given in Exp.\ II. Unfortunately, at the present time there is still no statement, even a heuristic one, that would encompass these two theorems (the first of which, for the moment, remains conjectural). A new idea is apparently lacking, as one is also lacking in Algebraic Geometry to obtain a proof of the Riemann--Roch theorem beyond projectivity hypotheses (cf.\ Exp.\ XIV, n\textsuperscript{o}~2).

We shall make no use in this seminar of the \underline{local theory} of intersections, set out in J.P. Serre, Algèbre Locale. Multiplicités (Lecture Notes in Mathematics n\textsuperscript{o}~11, Springer), and shall simply content ourselves with pointing out here that this course of Serre had an evident influence on the genesis of the theory in 1957. Nor shall we use the theory of the \underline{Chow ring} developed in the Chevalley Seminar 1958, Exposés 2 and 3 (École Normale Supérieure). This theory is essentially linked to nonsingularity hypotheses, whereas the aim of our seminar, on the contrary, is to develop an intersection theory on general schemes (and even on general locally ringed topoi) see on this subject the comments in Exp.\ XIV n\textsuperscript{o}~4 and n\textsuperscript{o}~8, giving the relations between our theory and that of the Chow ring, and raising the question of a generalization of the latter. We may say that the Seminar presents a coherent and self-contained intersection theory, but by no means exhausts the different points of view that are useful (or indeed indispensable) in intersection theory, and that it should consequently be supplemented by the cited exposés of Serre and Chevalley.

We have attached to the seminar (as an Appendix to Exp.\ 0) Grothendieck's 1957 report ``classes de Faisceaux et théorème de Riemann--Roch'', which served as the basis for the seminar of Borel and Serre at Princeton in the same year, as well as for their already cited article. This report outlines two proofs of the Riemann--Roch theorem for nonsingular quasi-projective varieties; the first, valid at present only in characteristic zero, but giving in return a slightly more precise result in the case of an immersion, does not yet appear in the literature (including the work of Borel--Serre and the present seminar). Reading this report of course presupposes none of the other exposés of the seminar, and may even serve as an introduction to its study on the same footing as Exposé 0, especially for a reader not yet familiar with Borel--Serre. Moreover, the proof to which we have just alluded applies essentially as it stands to the case
